\documentclass[10pt,a4paper]{amsart}
\usepackage[margin=19mm]{geometry}
\usepackage{amsmath,amssymb,mathtools}
\usepackage[scr=rsfs,cal=euler]{mathalfa}
\usepackage{xfrac}
\usepackage[unicode,hidelinks,bookmarks=false]{hyperref}
\newcommand{\ie}{i.e.\ }
\newcommand{\cf}{cf.\ }
\newcommand{\resp}{respectively\ }
\newcommand{\Subsection}{\subsection}
\providecommand{\nobreakdash}{\nobreak\hbox{-}}
\setlength{\emergencystretch}{2em}
\multlinegap0pt
\usepackage[shortlabels]{enumitem}
\usepackage{amssymb}
\usepackage{xcolor}
\usepackage{csquotes}
\newtheorem{prop}{Proposition}
\newtheorem{thm}{Theorem}
\newtheorem{lem}{Lemma}
\renewcommand{\d}{\mathrm{d}}
\newcommand{\ddc}{\d\d^c}
\def\tr{\operatorname{tr}}
\title{Determinantal capacity and $L^\infty$ estimates}
\author{Hao Fang, Biao Ma, Jinyang Wu}
\date{Source: arXiv:2609.24078v1 (CC BY 4.0)}
\begin{document}
\maketitle

\noindent\emph{Source and licence.} Determinantal capacity and $L^\infty$ estimates. Version arXiv:2609.24078v1 (CC BY 4.0), distributed by arXiv under the Creative Commons Attribution 4.0 International licence, http://creativecommons.org/licenses/by/4.0/, as recorded in the arXiv licence metadata for this exact version and shown on the abstract page https://arxiv.org/abs/2609.24078v1. Full licence notice: \texttt{licenses/arxiv-2609.24078-CC-BY-4.0.txt}.\par\smallskip

\noindent\textit{Editorial scope.} Counted unit: the article's prop prop:J-relative (source0.tex lines 3267--3305). The article's complete original proof of it is delivered with it (source0.tex lines 3307--3402, one \texttt{proof} environment of 96 lines). No mathematics is written, repaired or re-derived: the statement and the proof are the article's own lines, in the article's own notation, and no retained line is substituted or re-flowed.  Retained uncounted as the source-local context the proof needs: lines 482--491; lines 845--848; lines 858--871; lines 945--983; lines 1077--1086.  Nothing was omitted from any retained span, so the package carries no omission record.  No retained line contains a citation, so the wrapper carries no bibliography and no bibliography entry.  No retained line contains a figure environment, an image-inclusion command or drawing code, so no image is delivered and the package declares no asset dependency.  Editorial formatting only, in addition: an AMS \texttt{amsart} preamble that restates the article's own declarations and macros used by the retained lines, namely the environments \texttt{prop}, \texttt{thm}, \texttt{lem} on separate counters, together with the standard packages those lines need.  The retained environments print in this document as Proposition 1, Theorem 1, Lemma 1 in order of appearance, whereas the article prints the same items with the numbering of its own full document.  The complete native source of the article as distributed in the arXiv e-print for this exact version is bundled unchanged as \texttt{sources/211167/source0.tex}.
\begin{equation}
	\label{eq:main-equation}
	\left\{
	\begin{aligned}
		F_x\bigl(\omega_u(x)\bigr)&=\psi(x),\qquad x\in X,\\
		\omega_u(x)&=\omega(x)+\ddc u(x)\in\mathcal A_x,\\
		\sup_Xu&=0.
	\end{aligned}
	\right.
\end{equation}

	\begin{equation}
		\label{eq:viscosity-determinant-bound}
		\det_{\rho(x),+}\bigl(\chi(x)+\ddc q(x)\bigr)\leq e^{\gamma(x)}.
	\end{equation}

\begin{prop}
	\label{prop:determinant-reduction}
	Let $M\subset X$ be open, fix $\chi$ as above and a finite $\gamma\in C(M,\mathbb R)$, and let
	$\overline u$ be a viscosity supersolution of
	\eqref{eq:main-equation} on $M$. If $\underline u$ is a locally
	regularizable $\mathcal D$-subsolution, then
	$U=\overline u-\underline u$ satisfies the viscosity determinant
	bound on $M$.
	
	If $\overline u\in C^2(M)$ is an admissible supersolution, the same
	conclusion holds for every continuous viscosity
	$\mathcal D$-subsolution $\underline u$, without assuming local
	regularizability.
\end{prop}

\begin{thm}
	\label{thm:determinant-relative-Linfty}
	Let $(X, \rho)$ be a compact K\"{a}hler manifold, let
	$\chi$ be a smooth closed real $(1,1)$-form with $[\chi]$ big,
	and let $Z\subset X$ be closed and proper. Suppose that
	$\gamma\in C(X\setminus Z,[-\infty,+\infty))$ in the extended sense,
	that $e^\gamma>0$ almost everywhere on $X\setminus Z$, and that
	its zero extension belongs to $L^{p_0}(X,\rho^n)$ for some $p_0>1$.
	Let $U\in C(X,\mathbb R\cup\{+\infty\})$ satisfy
	$\{U=+\infty\}=Z$ and the viscosity determinant bound
	\eqref{eq:viscosity-determinant-bound} on $X\setminus Z$.
	Assume that, for some $M>0$,
	\begin{equation}
		\label{eq:thm-L1-hypothesis}
		\int_X(W_\chi-U)_+\,\rho^n\leq M,
	\end{equation}
	and that one of the following cases holds:
	\begin{enumerate}[label=\textnormal{(\alph*)}]
		\item\label{enu:thm-exp}
		\emph{Exponential case.} For some $\alpha>0$ and $q>n$,
		\[
		\operatorname{Ent}_{\mathcal D,q}(\chi,\gamma,\rho)\leq M,
		\qquad
		\sup_{\varphi\in\operatorname{PSH}_0(X,\chi)}
		\int_Xe^{-\alpha\varphi}\rho^n\leq M.
		\]
		\item\label{enu:thm-poly}
		\emph{Power case.} For some $p>1$ and $q>np'$, where $p'=p/(p-1)$,
		\[
		\int_X e^{p(\gamma-\log V_\chi)}\,\rho^n\leq M,
		\qquad
		\sup_{\varphi\in\operatorname{PSH}_0(X,\chi)}
		\int_X(-\varphi)^q\rho^n\leq M.
		\]
	\end{enumerate}
	Then $\sup_X(W_\chi-U)\leq C$. One may take
	$C=C(n,q,\alpha,M)$ in case~\ref{enu:thm-exp}, and
	$C=C(p,q,n,M,\int_X \rho^n)$ in case~\ref{enu:thm-poly}.
\end{thm}

\begin{lem}
	\label{lem:uniform-alpha-integrability}
	If $\chi\leq K\rho$ for some $K>0$, there are constants
	$\alpha,C_\alpha>0$, depending only on $(X,\rho,K)$, such that
	\begin{equation}
		\label{eq:uniform-alpha-integrability}
		\sup_{\varphi\in\operatorname{PSH}_0(X,\chi)}
		\int_Xe^{-\alpha\varphi}\rho^n\leq C_\alpha.
	\end{equation}
\end{lem}

\begin{prop}
\label{prop:J-relative}
Let $(X,\rho)$ be compact K\"ahler of dimension $n\geq 2$, let
$\omega,\alpha$ be K\"ahler forms, and let $\chi$ be a smooth closed real
$(1,1)$-form with $[\chi]$ big. For $A>0$, write
$$
P_\alpha(A)=\mu_1+\cdots+\mu_{n-1},
$$
where $\mu_1\geq\cdots\geq\mu_n>0$ are the eigenvalues of $\alpha$
relative to $A$.

Let $\underline u\leq 0$ be continuous on $X\setminus Z$, with
$\underline u\to-\infty$ along a proper closed set $Z$. Assume that,
whenever $q$ touches $\underline u$ from above,
$$
B_q:=\omega+\ddc q-\chi\geq a\alpha,
\qquad
P_\alpha(B_q)\leq c-\eta
$$
for fixed $a,\eta>0$ and $c>0$.

Let $s_j,r_j\geq 0$ be uniformly bounded and let
$u_j\in C^\infty(X)$, $\sup_Xu_j=0$, solve
$$
\omega_j:=(1+s_j)\omega+\ddc u_j>0,
\qquad
\tr_{\omega_j}\alpha+r_j\det_{\omega_j}\alpha=c.
$$
Then there exists $C>0$, independent of $j$, such that
$$
u_j\geq \underline u+W_\chi-C
\qquad\text{on }X\setminus Z.
$$
In particular, if $\chi\geq 0$, then $W_\chi=0$ and
$$
u_j\geq \underline u-C
\qquad\text{on }X\setminus Z.
$$
\end{prop}

\begin{proof}
Set
$$
F_j(A):=
\left(\tr_A\alpha+r_j\det_A\alpha\right)^{-1},
\qquad A>0.
$$
Then $F_j(\omega_j)=1/c$.

We first claim that $\underline u$ is a viscosity $\mathcal D$-subsolution of the
$j$-th equation, with shift $\chi$ and a determinant bound independent
of $j$. Let $q$ be an upper test for $\underline u$ and let $P\geq 0$
satisfy
$$
A:=(1+s_j)\omega+\ddc q-\chi+P>0,
\qquad
F_j(A)\leq \frac1c.
$$
Since
$$
A=B_q+s_j\omega+P\geq B_q\geq a\alpha,
$$
if $\mu_1\geq\cdots\geq\mu_n>0$ are the eigenvalues of $\alpha$
relative to $A$, then
$$
\mu_1+\cdots+\mu_{n-1}
=
P_\alpha(A)
\leq
P_\alpha(B_q)
\leq c-\eta.
$$
Moreover,
$$
\mu_1+\cdots+\mu_n+r_j\mu_1\cdots\mu_n\geq c.
$$
Hence, writing $r_*:=\sup_j r_j$ and using $\mu_k\leq a^{-1}$,
$$
\eta
\leq
\mu_n\left(1+r_*a^{-(n-1)}\right).
$$
Thus $\mu_n\geq\delta>0$ uniformly in $j$, so $A\leq\delta^{-1}\alpha$.
Since
$$
0\leq P\leq A,
$$
we obtain
$$
\det_\rho P\leq C_{\det}
$$
for a constant $C_{\det}$ independent of $j$. Therefore
$(\underline u,\chi,\gamma)$ is a viscosity $\mathcal D$-subsolution for the
$j$-th equation with
$$
\gamma:=\log C_{\det}.
$$

Set
$$
U_j:=u_j-\underline u
$$
on $X\setminus Z$, and extend $U_j$ by $+\infty$ on $Z$.
Since $u_j$ is smooth, Proposition~\ref{prop:determinant-reduction} gives
$$
\det_{\rho,+}(\chi+\ddc q)\leq e^\gamma
$$
for every lower test $q$ of $U_j$.

We now apply Theorem~\ref{thm:determinant-relative-Linfty}, case~\ref{enu:thm-exp}. Since $\gamma$ is constant,
its entropy is uniformly bounded. Also,
$$
(W_\chi-U_j)_+
=
(W_\chi+\underline u-u_j)_+
\leq -u_j.
$$
If $s_*:=\sup_j s_j$, then $\omega_j>0$ gives
$$
\ddc u_j\geq -(1+s_*)\omega.
$$
Together with $\sup_Xu_j=0$, the Green-function estimate yields
$$
\sup_j\int_X(-u_j)\rho^n<\infty.
$$
Finally, $\chi$ is fixed and smooth, so Lemma~\ref{lem:uniform-alpha-integrability} gives the required
uniform exponential integrability. Hence Theorem~\ref{thm:determinant-relative-Linfty} gives
$$
\sup_X(W_\chi-U_j)\leq C,
$$
with $C$ independent of $j$. This is exactly
$$
u_j\geq \underline u+W_\chi-C
$$
on $X\setminus Z$. If $\chi\geq 0$, then $W_\chi=0$.
\end{proof}

\end{document}
